%=============================================================================!
% 4.S  WHAT WE NOW HAVE                                                       !
%=============================================================================!
\unnumbered{What we now have}

A harmonic oscillator, $x = A\cos(\omega t + \phi)$, moves as the shadow of
a point going round a circle of radius $A$ at the angular speed $\omega$,
with its phase the starting angle on that circle. Its energy, $\frac12
kA^2$, is traded between kinetic and potential forms, which over a period
hold half each, and in the phase plane its states trace ellipses, circles
when the velocity is divided by $\omega$.

Near a smooth minimum with positive curvature, small swings are
approximately those of a spring of stiffness $U''$, which gave
the periods of a pendulum, of the bead of Chapter~3 and of a lithium atom
in a model hollow. The period of a swing of any size follows from its
energy as $2\int\dd x/v$, and in the cosine, Morse and Lennard-Jones wells
it grows with the energy, so small swings are the fastest.

A drag
proportional to the velocity makes the swing shrink by $\mathrm{e}^{-\gamma
t/2}$, or, if heavy enough, prevents any swing; a sinusoidal push gives a
steady response that peaks at resonance, with a height set by the quality
factor $\omega_0/\gamma$.

Two bodies on a spring oscillate in their
separation with the reduced mass $m_1m_2/(m_1 + m_2)$, and many coupled
bodies in the harmonic approximation move as sums of normal modes, whose patterns come from the
eigenvectors of the mass-weighted Hessian and whose squared angular
frequencies are its eigenvalues. The measured vibrations of
\cref{tab:os-vibrations} repeat every 9 to \SI{34}{\femto\second}, fastest
for bonds to hydrogen, and the fastest vibration present sets how short
the steps of a simulation must be.

Chapter~5 adds rotation, and with it the angular momentum that the forces
between atoms conserve.
